\hypertarget{appendix-c83016554b1e}{}\label{referencestart:appendix-c83016554b1e}
\section{Confidence intervals and sequential evidence}\label{appendix-c83016554b1e-confidence-intervals-and-sequential-evidence}

\textbf{Curriculum depth:} Developed · \textbf{First seen:} Phase 10 · \textbf{Latest development:} Phase 10

\begin{allmightroledefinition}

\subsection{Definition and mechanism}\label{appendix-c83016554b1e-definition-and-mechanism}

A confidence interval is a procedure with repeated-sample coverage under stated assumptions; it is not a posterior probability assigned to a fixed parameter. A normal interval for a mean uses a known population standard deviation or an approximation justified by the central limit theorem. With unknown scale and IID normal observations, the Student t distribution accounts for estimating the standard deviation. Heavy tails, sample dependence, multiple comparisons, and optional stopping can invalidate simple textbook guarantees. Adaptive algorithms that stop when a fixed-sample interval looks narrow need sequentially valid methods or a separately justified stopping rule. Report the target, sampling unit, dependence structure, estimator bias, interval procedure, stopping policy, and independent-run count. A narrow interval around a biased estimator measures precision about its expectation, not correctness of the target.

\end{allmightroledefinition}

\subsection{Prerequisites}\label{appendix-c83016554b1e-prerequisites}

\begin{itemize}
\tightlist
\item
  \hyperref[appendix-a2360b435383]{Common probability distributions and their assumptions}
\item
  \href{../05-Spatial-Worlds-Engine-State-and-Simulation/main.pdf\#nameddest=appendix-54c383df9992}{Expectation, covariance, and dependence (Spatial Worlds, Engine State, and Simulation, Reference notes: Expectation, covariance, and dependence)}
\item
  \href{../02-Mathematics-for-Visual-Computing/main.pdf\#nameddest=appendix-031e02bf28af}{Monte Carlo rendering estimators and variance (Mathematics for Visual Computing, Reference notes: Monte Carlo rendering estimators and variance)}
\end{itemize}

\subsection{Used by}\label{appendix-c83016554b1e-used-by}

\begin{itemize}
\item
  This node is a current endpoint in the supplied Phase 9--10 path.
\item
  \hyperref[appendix-b9db93e28371]{Reproducible numerical benchmarking} \#\# Full teaching explanation
\item
  \href{../02-Mathematics-for-Visual-Computing/main.pdf\#nameddest=02-foundations-ii-ch030-inference-and-error-metrics-need-explicit-assumptions}{Phase 10 teaching section (Mathematics for Visual Computing, Ch. 14)}
\item
  \href{../01-Bits-Programs-and-Reliable-Systems/main.pdf\#nameddest=appendix-cb5af353ef77}{Visual Recall/10 - Probability and Monte Carlo - Visual Recall.md (Bits, Programs, and Reliable Systems, Reference notes: 10 - Probability and Monte Carlo - Visual Recall)} · visual recall
\end{itemize}

\subsection{Source coverage}\label{appendix-c83016554b1e-source-coverage}

Supplied lesson parts: 960 (T06 · M058), 961 (T06 · M058).
